\documentclass[12pt]{amsart}


\usepackage{times}
\usepackage[margin=.71 in]{geometry}
\usepackage{paralist,amsmath,amssymb,multicol,graphicx,framed,ifthen,color,xcolor,stmaryrd,enumitem,colonequals}
\usepackage[outline]{contour}
\contourlength{.4pt}
\contournumber{10}
\newcommand{\Bold}[1]{\contour{black}{#1}}

\definecolor{chianti}{rgb}{0.6,0,0}
\definecolor{meretale}{rgb}{0,0,.6}
\definecolor{leaf}{rgb}{0,.35,0}
\newcommand{\Q}{\mathbb{Q}}
\newcommand{\N}{\mathbb{N}}
\newcommand{\Z}{\mathbb{Z}}
\newcommand{\R}{\mathbb{R}}
\newcommand{\C}{\mathbb{C}}
\newcommand{\e}{\varepsilon}
\newcommand{\inv}{^{-1}}
\newcommand{\dabs}[1]{\left| #1 \right|}
\newcommand{\ds}{\displaystyle}
\newcommand{\solution}[1]{\ifthenelse {\equal{\displaysol}{1}} {\begin{framed}{\color{meretale}\noindent #1}\end{framed}} { \ }}
\newcommand{\solutione}[1]{\ifthenelse {\equal{\displaysol}{1}} {\begin{framed}{\color{leaf}This solution is embargoed.}\end{framed}} { \ }}
\newcommand{\showsol}[1]{\def\displaysol{#1}}
\newcommand{\nosolfill}{\ifthenelse {\equal{\displaysol}{1}} {} {\vfill}}
\newcommand{\nosolpage}{\ifthenelse {\equal{\displaysol}{1}} {} {\newpage}}
\newcommand{\nosolonly}[1]{\ifthenelse {\equal{\displaysol}{1}} {} {{#1}}}

\newcommand{\rsa}{\rightsquigarrow}


\newcommand\itemA{\stepcounter{enumi}\item[{\Bold{(\theenumi)}}]}
\newcommand\itemB{\stepcounter{enumi}\item[(\theenumi)]}
\newcommand\itemC{\stepcounter{enumi}\item[{\it{(\theenumi)}}]}
\newcommand\itema{\stepcounter{enumii}\item[{\Bold{(\theenumii)}}]}
\newcommand\itemb{\stepcounter{enumii}\item[(\theenumii)]}
\newcommand\itemc{\stepcounter{enumii}\item[{\it{(\theenumii)}}]}
\newcommand\itemai{\stepcounter{enumiii}\item[{\Bold{(\theenumiii)}}]}
\newcommand\itembi{\stepcounter{enumiii}\item[(\theenumiii)]}
\newcommand\itemci{\stepcounter{enumiii}\item[{\it{(\theenumiii)}}]}
\newcommand\ceq{\colonequals}


\DeclareMathOperator{\ord}{ord}

\DeclareMathOperator{\res}{res}
\setlength\parindent{0pt}
%\usepackage{times}

%\addtolength{\textwidth}{100pt}
%\addtolength{\evensidemargin}{-45pt}\mathrm{Gal}
%\addtolength{\oddsidemargin}{-60pt}

\pagestyle{empty}
%\begin{document}\begin{itemize}

%\thispagestyle{empty}




\begin{document}
\showsol{0}
	
	\thispagestyle{empty}
	
	\section*{The Fundamental Theorem of Galois Theory}
	\begin{framed}
	\textsc{The Fundamental Theorem of Galois Theory:} 
	Let $L/F$ be a finite Galois extension. There is a bijection 
	\[ \begin{array}{ccccc} & \{ \text{fields} \ E \ | \ F\subseteq E \subseteq L \} &\longleftrightarrow & \{ \text{subgroups} \ H \leq \mathrm{Gal}(L/F) \} & \\
	& E & {\mapsto} & \mathrm{Gal}(L/E) & \\
	& L^H & {\mapsfrom} & H &
	\end{array}\]

Moreover, this bijective correspondence enjoys the following properties:
\begin{enumerate}
\item $\stackrel{\mathrm{Gal}(L/\star)}{\longrightarrow}$ and $\stackrel{L^{\star}}{\longleftarrow}$ each reverse the order of inclusion: 
\begin{itemize}
\item[] $E \subseteq E' \Leftrightarrow \mathrm{Gal}(L/E') \leq \mathrm{Gal}(L/E)$
\end{itemize}
\item $\stackrel{\mathrm{Gal}(L/\star)}{\longrightarrow}$ and $\stackrel{L^{\star}}{\longleftarrow}$ convert between degrees of extensions and indices of subgroups:  
\begin{itemize}
\item[] $[E:F] = [\mathrm{Gal}(L/F) : \mathrm{Gal}(L/E)]$; \ equivalently
\item[] $[\mathrm{Gal}(L/F) : H] = [L^H:F]$.
\end{itemize}
\item Normal subgroups correspond to intermediate fields that are Galois over $F$: 
\begin{itemize}
\item[] $E/F$ is Galois if and only if $\mathrm{Gal}(L/E) \trianglelefteq \mathrm{Gal}(L/F)$; \ equivalently
\item[] $N \trianglelefteq G$ if and only if $L^N/F$ is Galois.
\end{itemize}
In this case, $\mathrm{Gal}(E/F) = \mathrm{Gal}(L/F) / \mathrm{Gal}(L/E)$.
\end{enumerate}
		\end{framed}

\

	
	\begin{enumerate}
	\itemA Let $L$ be the splitting field of $x^3-2$ over $\Q$. From earlier work, we know that:
\medskip
	\begin{itemize}
	\item $L = \Q(\sqrt[3]{2},\omega)$ where $\omega=e^{2\pi i/3}$. 
	\smallskip
	\item $\mathrm{Gal}(L/\Q) =\langle \sigma, \tau \rangle$, with $\sigma\begin{bmatrix} \sqrt[3]{2} \\ \omega \end{bmatrix} =\begin{bmatrix} \omega \sqrt[3]{2} \\ \omega \end{bmatrix}$ and $\tau \begin{bmatrix} \sqrt[3]{2} \\ \omega \end{bmatrix} =\begin{bmatrix} \sqrt[3]{2} \\ \omega^2 \end{bmatrix}$.
		\smallskip
	\item $\mathrm{Gal}(L/\Q) \cong S_3$ via $\sigma \mapsto (1\ 2 \ 3)$ and $\tau \mapsto (2\ 3)$.
	\end{itemize}
\medskip
	\begin{enumerate}
	\itema List all subgroups of $S_3$, and draw a diagram indicating which is contained in which.
	\itema Use the previous part and the Fundamental Theorem to determine how many\footnote{Do \emph{not} explicitly identify the fields yet; just make up a placeholder name for each intermediate field $E$.} intermediate fields $E$ are between $\Q$ and $L$. Then draw a diagram to indicate containments among the fields.
	\itema Compute the index of each subgroup of $S_3$. Use this and the Fundamental Theorem to determine $[E:\Q]$ for each intermediate field $E$.
	\itema Determine which subgroups of $S_3$ are normal. Use this and the Fundamental Theorem to determine which intermediate fields are Galois over $\Q$.
	\itema Explicitly identify the intermediate fields $E$.
	\end{enumerate}
	
		\
		
		\itemA Show that the maps $\stackrel{\mathrm{Gal}(L/\star)}{\longrightarrow}$ and $\stackrel{L^{\star}}{\longleftarrow}$ give the bijection as stated in the Theorem.
		
		\
	
	\itemB Let $L/F$ be a finite Galois extension.
	\begin{enumerate}
	\itemb When is every intermediate field $E\subseteq L$ Galois over $F$?
	\itemb If $[L:F]=n$ and $n=p^e m$ for some prime $p$ and positive integer $m$ that is not a multiple of $p$, explain why there exists an intermediate field $E$ with $[E:F]=p^e$.
	\end{enumerate}
	
	\
	
	\itemB Show that if $L/F$ is Galois with $[L:F]=55$, then there is a unique intermediate field $E$ with $[E:F]=5$.
	
	\
	
	\itemB Let $L$ be the splitting field of $x^4-2026$ over $\Q$. Show that there exists a unique field $E\subseteq L$ such that $[E:\Q]=4$ and $E/\Q$ is Galois.
	
	\end{enumerate}
	

	
	\newpage
	



\end{document}
